\documentclass[11pt,a4paper]{article}

% --- Packages ---
\usepackage[margin=1in]{geometry}
\usepackage{amsmath,amssymb,amsfonts}
\usepackage{booktabs}
\usepackage{tabularx}
\usepackage{multirow}
\usepackage{enumitem}
\usepackage{xcolor}
\usepackage{titlesec}
\usepackage{caption}
\usepackage{cite}
\usepackage[hidelinks,colorlinks=true,linkcolor=blue,citecolor=blue,urlcolor=blue]{hyperref}

% --- Custom Styling ---
\definecolor{primary}{RGB}{30, 58, 138}    % Deep Navy
\definecolor{secondary}{RGB}{15, 118, 110} % Dark Teal
\definecolor{lightbg}{RGB}{243, 244, 246}  % Soft Light Gray

\titleformat{\section}{\large\bfseries\color{primary}}{\thesection.}{0.5em}{}[\titlerule]
\titleformat{\subsection}{\normalfont\bfseries\color{secondary}}{\thesubsection}{0.5em}{}

\setlist[itemize]{noitemsep, topsep=3pt, leftmargin=1.5em}
\setlist[enumerate]{noitemsep, topsep=3pt, leftmargin=1.5em}

% --- Document Metadata ---
\title{\textbf{\Large Coding Agent Adoption in JavaScript / TypeScript / npm\\}\\[0.5em]
\Large\textbf{Thesis Proposal}}

\author{Manuel Youssef Haik Kevorkian}
\date{October 2026}

\begin{document}

\maketitle

% ─────────────────────────────────────────────────────────────────
\section{Aim \& Motivation}
% ─────────────────────────────────────────────────────────────────

\subsection{Background}

Autonomous AI coding agents such as Claude Code, Cursor, GitHub Copilot, Windsurf, Aider,
and OpenAI Codex CLI have fundamentally changed how software is written on GitHub. Unlike
first-generation autocomplete assistants, these tools plan tasks, write multi-file code, run
tests, and open pull requests with minimal human input. Recent large-scale studies estimate
adoption rates of 22--28\% across all established GitHub repositories \cite{robbes2026agenticmuch}
and 71--76\% in newly created repositories \cite{robbes2026agenticverymuch}, confirming that
agentic development has moved from experimental to mainstream.

Despite this rapid growth, two significant gaps remain in the literature. First, no study
has measured adoption specifically within the JavaScript/TypeScript/npm ecosystem, even
though configuration file evidence and GitHub Octoverse data show TypeScript to be the
most active ecosystem for agent adoption \cite{galster2026dataset, github2025octoverse}.
Second, geographic variation in AI coding tool adoption has been documented at a global
level \cite{daniotti2025} but never studied within a specific ecosystem with a multi-signal
detection pipeline and formal statistical tests at the country level.

\subsection{Research Gaps and Questions}

\textbf{Gap 1 (Ecosystem scope):} All major adoption studies average across all GitHub
languages, masking ecosystem-specific patterns. The JS/TS/npm ecosystem, home to the
world's largest package registry, has never been studied in isolation.

\textbf{Gap 2 (Undercounting):} Single-signal detection misses up to 30$\times$ the
AI usage captured by multi-signal approaches \cite{khosravani2026detecting}. No study
has quantified this undercounting delta for the JS/TS ecosystem, where npm SDK dependency
declarations provide an additional signal family unavailable in other languages.

\textbf{Gap 3 (Geographic ranking):} Daniotti et al.~\cite{daniotti2025} document
country-level variation in AI adoption but focus on Python, use function-level detection
only, and do not apply formal statistical tests. No study has ranked countries by adoption
rate within the JS/TS/npm ecosystem with confidence intervals and significance testing.

\medskip
\noindent These gaps define three research questions:

\begin{itemize}
    \item \textbf{RQ1 (Adoption scope):} What proportion of active JS/TS/npm repositories
          on GitHub show evidence of AI coding agent use, and is this proportion
          statistically non-trivial?
    \item \textbf{RQ2 (Undercounting):} How much does single-signal detection undercount
          AI adoption compared to a multi-signal approach, and is this difference
          statistically significant?
    \item \textbf{RQ3 (Geographic ranking):} Do AI coding agent adoption rates differ
          significantly across countries within the JS/TS/npm ecosystem, and which
          countries show the highest and lowest rates?
\end{itemize}

% ─────────────────────────────────────────────────────────────────
\section{Methodology}
% ─────────────────────────────────────────────────────────────────

The study follows a four-phase computational pipeline implemented in Python using the
GitHub GraphQL and REST APIs. The full corpus of \textbf{60,129 repositories} is processed
end-to-end with no sampling, because the layered detection design classifies most repositories
cheaply in early passes, making sampling unnecessary and avoiding bias in the country-level analysis.

\subsection{Phase 1: Corpus Construction}

Repositories were collected via the GitHub GraphQL Search API using rolling 5-day date windows
from 2025-02-24 to October 2026, with recursive window halving when the 1,000-result cap was
hit. Inclusion filters (agreed across the research team) are shown below.

\begin{center}
\begin{tabular}{@{}lll@{}}
\toprule
\textbf{Filter} & \textbf{Value} & \textbf{Rationale} \\
\midrule
Language            & JavaScript or TypeScript & Thesis scope \\
Created on or after & 2025-02-24               & Broad agentic-era availability \\
Stars               & $\geq 4$                 & Excludes toy / abandoned repos \\
Commits             & $\geq 50$                & Meaningful development history \\
Fork                & No                       & Original projects only \\
JS + TS bytes       & $\geq 175{,}000$         & JS/TS is primary language \\
\bottomrule
\end{tabular}
\end{center}

Each repository is categorized by package manager (npm, pnpm, yarn, bun, deno, vanilla)
using lockfile presence at the root. Result: \textbf{60,129 repositories} spanning
all JS/TS ecosystems.

\subsection{Phase 2: Owner Location Inference}

GitHub location is attached to user profiles, not repositories. For the 50,002 unique
owner logins in the corpus, country is inferred in two passes:

\begin{itemize}
    \item \textbf{Pass 1} (fast): Batch GraphQL queries (50 owners/request) match the raw
          \texttt{location} field against a keyword table of country names and major cities.
          Result: 13,625 owners resolved (27.2\%).
    \item \textbf{Pass 2} (deep): One query per unresolved owner tries six signals in priority
          order: bio keyword, company keyword, company domain TLD, email domain TLD,
          website TLD, organisation location, and network majority vote.
          Owners unresolved after all signals are labelled \texttt{not\_specified} and
          excluded from country-level analysis but retained in global totals.
\end{itemize}

\subsection{Phase 3: AI Adoption Detection (3-Pass Layered Pipeline)}

Passes are ordered by API cost. A repository classified in an early pass is never
re-examined, ensuring the most expensive checks run on the smallest possible subset.

\begin{center}
\begin{tabularx}{\textwidth}{@{}p{1.6cm}Xp{3.4cm}@{}}
\toprule
\textbf{Pass} & \textbf{What is checked} & \textbf{Outcome if found} \\
\midrule
Pass 1 & 12 known agent config files via batch GraphQL (10 repos/request):
         \texttt{CLAUDE.md}, \texttt{.cursorrules}, \texttt{GEMINI.md},
         \texttt{.aider.conf.yml}, \texttt{copilot-instructions.md}, etc. &
         Immediately labelled \textbf{agentic}; skips Pass 2 and 3 \\
\addlinespace
Pass 2 & 100 most recent commits + 50 most recent PRs per repo.
         Scans \texttt{Co-authored-by:} trailer patterns and author email patterns.
         PR branch names scanned for \texttt{copilot/}, \texttt{devin/}, \texttt{codex/} prefixes. &
         Contributes signals to final label \\
\addlinespace
Pass 3 & Up to 300 commits + 100 PRs (body, labels, branch) + \texttt{package.json}
         dependency file. Checks npm SDK deps: \texttt{@anthropic-ai/sdk},
         \texttt{@anthropic-ai/claude-code}, \texttt{@cursor/agent}. &
         Contributes signals to final label \\
\bottomrule
\end{tabularx}
\end{center}

\medskip
\noindent\textbf{Classification rule:}
$\geq$3 signals OR $\geq$2 distinct agents $\Rightarrow$ \textbf{agentic};
\quad 1--2 signals from a single agent $\Rightarrow$ \textbf{ai\_assisted};
\quad no signal $\Rightarrow$ \textbf{none}.

Pass 1 config-file-only labels are saved separately alongside the full three-pass labels
to directly answer RQ2 (undercounting delta).

\subsection{Phase 4: Geographic Analysis and Statistical Tests}

The detected corpus is joined with owner locations on \texttt{owner\_login}. For each
country with $\geq$30 repositories, adoption rate (\%) and Wilson 95\% confidence interval
are computed. Countries are ranked by adoption rate (descending).

\begin{itemize}
    \item \textbf{RQ1}: One-sample binomial test, null: global adoption rate $\leq$ 5\%.
    \item \textbf{RQ2}: McNemar's test on paired per-repo labels (Pass 1 only vs.\ full pipeline).
    \item \textbf{RQ3}: Chi-squared test of independence across all eligible countries,
          followed by pairwise Fisher's exact tests with Benjamini-Hochberg FDR correction
          at $\alpha = 0.05$.
\end{itemize}

% ─────────────────────────────────────────────────────────────────
\section{Literature}
% ─────────────────────────────────────────────────────────────────

Thirty papers were gathered and reviewed. They fall into five thematic groups:

\begin{enumerate}[leftmargin=1.5em, itemsep=4pt]
    \item \textbf{Core adoption studies:} Robbes et al.\ (Agentic Much? \cite{robbes2026agenticmuch};
          Agentic Very Much!\ \cite{robbes2026agenticverymuch}); Khosravani \& Mockus
          \cite{khosravani2026detecting}. Establish global baselines (22--28\% established
          repos; 71--76\% new repos; 30$\times$ undercounting from single-signal detection).

    \item \textbf{Detection methodology:} Robbes et al.\ heuristics catalogue
          \cite{robbes2026promises}; Ghaleb agent fingerprinting \cite{ghaleb2026fingerprinting};
          Cynthia et al.\ usage patterns \cite{cynthia2026};
          Dey et al.\ BIMAN bot detection \cite{dey2020biman}.

    \item \textbf{Configuration artifacts and ecosystem:} Galster et al.\ configuration artifacts
          \cite{galster2026dataset}, configuring agents \cite{galster2026configuring},
          harness engineering \cite{galster2026harness};
          Chatlatanagulchai et al.\ context files \cite{chatlatanagulchai2025agent};
          Dakhel et al.\ Copilot quality \cite{dakhel2023};
          Chen et al.\ Codex \cite{chen2021codex};
          Agarwal et al.\ AI IDEs vs agents \cite{agarwal2026}.

    \item \textbf{PR dynamics and code quality:} Ogenrwot \& Businge \cite{ogenrwot2026how};
          Mazloomzadeh et al.\ \cite{mazloomzadeh2026how};
          Popescu et al.\ \cite{popescu2026investigating};
          Watanabe et al.\ \cite{watanabe2025use};
          Oukhay et al.\ \cite{oukhay2026when};
          Hora \& Robbes \cite{hora2026are};
          Barke et al.\ \cite{barke2023};
          Ziegler et al.\ \cite{ziegler2022}.

    \item \textbf{Research infrastructure and geography:} Kalliamvakou et al.\
          \cite{kalliamvakou2016promises}; Dabic et al.\ GHS \cite{dabic2021sampling};
          Chatzidimitriou et al.\ npm-miner \cite{chatzidimitriou2018};
          Wachs et al.\ geography of OSS \cite{wachs2022geography};
          Daniotti et al.\ global AI diffusion \cite{daniotti2025};
          Felder et al.\ German industry adoption \cite{felder2026};
          GitHub Octoverse 2025 \cite{github2025octoverse};
          GitHub Octoverse 2024 \cite{github2024octoverse}.
\end{enumerate}

% ─────────────────────────────────────────────────────────────────
\section{Expected Outcomes}
% ─────────────────────────────────────────────────────────────────

\begin{enumerate}[leftmargin=1.5em, itemsep=4pt]
    \item \textbf{First ecosystem-scoped adoption estimate} for the JS/TS/npm ecosystem,
          broken down by package manager (npm, pnpm, yarn, bun, deno, vanilla).
    \item \textbf{Quantified undercounting delta}: a measured percentage-point comparison
          between single-signal (config file only) and multi-signal (three-pass) detection
          rates on the same 60,129-repository corpus.
    \item \textbf{Country-level adoption ranking} with Wilson 95\% confidence intervals and
          significance-tested pairwise comparisons across all countries with $\geq$30
          repositories in the corpus.
\end{enumerate}

% ─────────────────────────────────────────────────────────────────
\section{Work Plan \& Timeline}
% ─────────────────────────────────────────────────────────────────

\begin{center}
\begin{tabularx}{\textwidth}{@{}l l X@{}}
\toprule
\textbf{Phase} & \textbf{Dates} & \textbf{Work} \\
\midrule
Phase 0 & Sep 12 -- Sep 20, 2026 &
    Background reading, gap identification, project plan. \\
\addlinespace
\textit{Milestone} & \textit{Sep 20, 2026} & \textit{Project plan / strategy submitted} \\
\addlinespace
Phase 1 & Sep 21 -- Oct 13, 2026 &
    Corpus collection (60,129 repos), ecosystem categorization, owner location
    fast pass (27.2\% resolved), config-file detection pass. \\
\addlinespace
\textit{Milestone} & \textit{Oct 13, 2026} & \textit{First interim report submitted} \\
\addlinespace
Phase 2 & Oct 14 -- Nov 15, 2026 &
    Owner location deep pass (in progress), commit trailer and deep scan passes,
    geographic join, statistical tests. Write Methodology and Results chapters. \\
\addlinespace
\textit{Milestone} & \textit{Nov 15, 2026} & \textit{Second interim report submitted} \\
\addlinespace
Phase 3 & Nov 16 -- Dec 10, 2026 &
    Write Discussion and Conclusion chapters. Compile full thesis draft. \\
\addlinespace
\textit{Milestone} & \textit{Dec 10, 2026} & \textit{Full thesis draft complete} \\
\addlinespace
Phase 4 & Dec 11, 2026 -- Jan 18, 2027 &
    Incorporate supervisor feedback; polish writing; prepare presentation. \\
\addlinespace
Phase 5 & Jan 19 -- Jan 27, 2027 &
    Oral presentation (20 min) and Q\&A with supervisor and reviewer. \\
\addlinespace
Phase 6 & Jan 28 -- Jan 31, 2027 &
    Incorporate reviewer comments; finalize. \\
\addlinespace
\textit{Milestone} & \textit{Jan 31, 2027} & \textit{Final submission: hard copy + digital copy} \\
\bottomrule
\end{tabularx}
\end{center}

\subsection{Threats to Validity}

\textbf{Detection recall.} Heuristic-based detection cannot guarantee full recall.
Reported adoption rates should be interpreted as a lower bound on true agent usage,
consistent with the treatment in Robbes et al.~\cite{robbes2026agenticmuch}.

\textbf{Location attribution accuracy.} Approximately 40\% of owners remain unresolved
after both location passes and are excluded from country-level analysis. If unresolved
owners are systematically concentrated in specific countries, adoption rates for those
countries may be underestimated. The two-pass design reduces this bias relative to a
single field-lookup approach \cite{wachs2022geography}.

\textbf{Heuristic currency.} Agent vendors may change their default commit patterns
after data collection. All detection signals are version-dated, and the study reports
the period of validity for each heuristic following the approach of
Robbes et al.~\cite{robbes2026promises}.

\newpage

\begin{thebibliography}{99}

\bibitem{galster2026dataset}
M. Galster et al.,
``Configuration Artifacts of AI Coding Agents: An Empirical Study,''
in \textit{Proc. ACM AIware '26}, 2026. DOI: 10.1145/3805760.3814922.

\bibitem{khosravani2026detecting}
A. Khosravani and A. Mockus,
``A Census of AI-Assisted Software Development on GitHub,''
\textit{arXiv:2606.24429}, 2026.

\bibitem{ghaleb2026fingerprinting}
T. A. Ghaleb,
``Fingerprinting AI Coding Agents on GitHub,''
in \textit{Proc. MSR '26 Mining Challenge}, 2026.
arXiv:2601.17406. DOI: 10.1145/3793302.3793582.

\bibitem{robbes2026promises}
R. Robbes et al.,
``Heuristics for Detecting AI-Generated Commits on GitHub,''
in \textit{Proc. MSR '26}, 2026. DOI: 10.1145/3793302.3793375.

\bibitem{robbes2026agenticmuch}
R. Robbes et al.,
``Agentic Much? Adoption of Coding Agents on GitHub,''
\textit{arXiv:2601.18341}, 2026.

\bibitem{robbes2026agenticverymuch}
R. Robbes et al.,
``Agentic Very Much! Adoption of Coding Agents in New GitHub Projects,''
\textit{arXiv:2606.07448}, 2026.

\bibitem{cynthia2026}
S. T. Cynthia, J. K. Das, and B. Roy,
``Are We All Using Agents the Same Way? An Empirical Study of Core and
Peripheral Developers Use of Coding Agents,''
\textit{arXiv:2601.20106}, 2026.

\bibitem{chatlatanagulchai2025agent}
W. Chatlatanagulchai et al.,
``Context Files for AI Coding Agents: Usage and Characteristics,''
\textit{arXiv:2511.12884}, 2025.

\bibitem{ogenrwot2026how}
E. Ogenrwot and J. Businge,
``How AI Coding Agents Modify Code: A Large-Scale Study of GitHub Pull Requests,''
in \textit{Proc. MSR '26 Mining Challenge}, 2026.
arXiv:2601.17581. DOI: 10.1145/3793302.3793603.

\bibitem{mazloomzadeh2026how}
I. Mazloomzadeh, M. M. Morovati, and F. Khomh,
``How Do AI Coding Agents Contribute to Software Development?
An Empirical Study of Agentic Pull Requests,''
\textit{arXiv:2607.21832}, 2026.

\bibitem{popescu2026investigating}
R. M. Popescu, D. Gros, A. Botocan, R. Pandita, P. Devanbu, and M. Izadi,
``Investigating Autonomous Agent Contributions in the Wild:
Activity Patterns and Code Change over Time,''
in \textit{Proc. MSR '26}, 2026. arXiv:2604.00917.

\bibitem{watanabe2025use}
M. Watanabe, H. Li, Y. Kashiwa, B. Reid, H. Iida, and A. E. Hassan,
``On the Use of Agentic Coding: An Empirical Study of Pull Requests on GitHub,''
\textit{arXiv:2509.14745}, 2025.

\bibitem{oukhay2026when}
H. Oukhay et al.,
``Do AI-Generated Commits Get Reverted More Often?''
in \textit{Proc. MSR '26}, 2026. DOI: 10.1145/3793302.3793587.

\bibitem{hora2026are}
A. Hora and R. Robbes,
``How AI Coding Agents Affect Mocking Practices in Test Suites,''
in \textit{Proc. MSR '26}, 2026. DOI: 10.1145/3793302.3793362.

\bibitem{galster2026configuring}
M. Galster et al.,
``Configuring Agentic AI Coding Tools: An Exploratory Study,''
in \textit{Proc. ACM AIware '26}, 2026. DOI: 10.1145/3805760.3814887.

\bibitem{galster2026harness}
M. Galster et al.,
``Harness Engineering for Agentic AI Coding Tools: An Exploratory Study,''
in \textit{Proc. ACM AIware '26}, 2026. arXiv:2602.14690.

\bibitem{dakhel2023}
A. M. Dakhel, V. Majdinasab, A. Nikanjam, F. Khomh, M. C. Desmarais,
and Z. M. Jiang,
``GitHub Copilot AI Pair Programmer: Asset or Liability?''
\textit{Journal of Systems and Software}, vol. 203, art. 111734, 2023.
DOI: 10.1016/j.jss.2023.111734.

\bibitem{dabic2021sampling}
O. Dabi\'{c} et al.,
``Sampling Projects in GitHub for MSR Studies,''
in \textit{Proc. MSR '21 Data Showcase}, 2021. DOI: 10.5281/zenodo.4588464.

\bibitem{wachs2022geography}
J. Wachs et al.,
``The Geography of Open Source Software: Evidence from GitHub,''
\textit{Technological Forecasting and Social Change}, vol. 176, art. 121478, 2022.

\bibitem{daniotti2025}
S. Daniotti, J. Wachs, X. Feng, and F. Neffke,
``Who is using AI to code? Global diffusion and impact of generative AI,''
\textit{arXiv:2506.08945}, 2025.

\bibitem{felder2026}
L. Felder et al.,
``Adoption of Generative AI in the German Software Engineering Industry,''
in \textit{Proc. FSE '26}, 2026. \textit{arXiv:2601.16700}.

\bibitem{github2025octoverse}
GitHub, Inc.,
``GitHub Octoverse 2025,'' GitHub Blog, 2025.

\bibitem{github2024octoverse}
GitHub, Inc.,
``GitHub Octoverse 2024,'' GitHub Blog, 2024.

\bibitem{kalliamvakou2016promises}
E. Kalliamvakou et al.,
``An In-Depth Study of the Promises and Perils of Mining GitHub,''
\textit{Empirical Software Engineering}, vol. 21, no. 5, pp. 2035--2071, 2016.

\bibitem{barke2023}
S. Barke, M. B. James, and N. Polikarpova,
``Grounded Copilot: How Programmers Interact with Code-Generating Models,''
\textit{Proc. ACM Program. Lang.}, vol. 7, OOPSLA1, art. 78, 2023.

\bibitem{ziegler2022}
A. Ziegler et al.,
``Productivity Assessment of Neural Code Completion,''
in \textit{Proc. MAPS '22}, pp. 21--29, 2022.

\bibitem{chatzidimitriou2018}
K. C. Chatzidimitriou et al.,
``npm-miner: An Infrastructure for Measuring the Quality of the npm Registry,''
in \textit{Proc. MSR '18}, pp. 42--45, 2018.

\bibitem{chen2021codex}
M. Chen et al.,
``Evaluating Large Language Models Trained on Code,''
\textit{arXiv:2107.03374}, 2021.

\bibitem{dey2020biman}
T. Dey et al.,
``Detecting and Characterizing Bots that Commit Code,''
in \textit{Proc. MSR '20}, 2020. \textit{arXiv:2003.03172}.

\bibitem{agarwal2026}
S. Agarwal, H. He, and B. Vasilescu,
``AI IDEs or Autonomous Agents? Measuring the Impact of Coding Agents on Software Development,''
\textit{arXiv:2601.13597}, 2026.

\end{thebibliography}

\end{document}
